\clearpage
\noindent\textbf{虚拟互动、心笺与失败恢复}\par\smallskip
\label{sec:impl-10-virtual}
\par\medskip\noindent\textbf{虚拟回复的两阶段写入。}
虚拟对话由\DesignCode{public/core.mjs}中的状态规则与服务端动作层协作完成。用户先提交文本，服务端保存用户消息和待回复任务，再调用模型生成对方回复。只有内部路径可以执行\DesignCode{finishReply}；外部伪造同名动作会被拒绝。

模型超时或网络失败时保留待回复状态。再次重试只补齐回复，不再次发送同一条用户消息。等待任务未处理完时，虚拟状态机拒绝继续写入冲突消息；这与真人聊天允许边分析边手动发送是不同的业务约束。

\begin{table}[H]\centering\small
\caption{虚拟状态与私人记录的分界}
\begin{tabularx}{\linewidth}{@{}p{30mm}X@{}}\toprule
状态或操作 & 实现含义 \\ \midrule
待回复 & 用户消息已保存，模型尚未成功补写 \\
心笺分数与备注 & 用户私人记录，不自动形成对方意愿 \\
回望与授权 & 当前周期的用户选择，跨周期需重新确认 \\
屏蔽 & 停止相关虚拟互动，并撤回对应授权 \\
演示重置 & 重建虚拟状态和默认资料，不清空真人消息与媒体 \\ \bottomrule
\end{tabularx}\end{table}

\par\medskip\noindent\textbf{时间、资格与重置。}
每七天形成一个虚拟周期，“轻叩窗扉”先检查双向交流、当期意愿及双方授权；资格不满足不消耗机会，具体分数不外发。测试通过固定时间覆盖跨期、拒绝和撤回，不使用真实等待七天的方式验证时间规则。

Demo注入来信、日期推进和重置属于演示入口，不表示真人已经执行了相同行为。日期推进仅接受合法整数并受累计范围约束。测试检查失败动作不写脏状态，以及上游素材不包含心笺备注等私人字段，避免把模型回复当作系统授权依据。
